\printedpage{199}
% head: A. Grothendieck : Dérivateurs, Ch. XVIII, Édition des Universités de Montpellier II et Paris VII
% foot: Version 17 décembre 2025

\noindent\textbf{Corollary.} \textit{Let $\pi$ be an infinite cardinal. For the closed subcategory $\underline{S}'$ of $\mathcal{M}$ (consisting of commutative $k$-algebras $A$ such that $\Hom_{\mathcal{M}}(A,k)=\emptyset$) to be $\pi$-accessible, it is necessary and sufficient that for every $k$-algebra $A$ (of cardinality ${}>\pi$), such that $\Hom_{\mathcal{M}}(A,k)=\emptyset$, there exist a subalgebra $A_i$} of cardinality $\leqslant\pi$, \textit{such that $\Hom_{\mathcal{M}}(A_i,k)=\emptyset$.} % [sic] tout k-algèbre -> toute k-algèbre, « algèbre » est féminin

\noindent(We apply the lemma, taking for $\alpha$ the first infinite ordinal.)

\noindent\textbf{Question.} \textit{Does this hold, simply taking $\pi\geqslant\card k$~?} The question of the existence of a cardinal $\pi$, having the property indicated above, seems delicate to me, and I do not know the answer. If one wants to construct a counterexample for every $\pi$, the preceding developments suggest taking for $A$ an infinite product of nontrivial extensions of $k$, or a ring $k_{\Phi}$ as in the proof of proposition 4 a).